\documentclass{article}
\usepackage[pdftex]{graphicx}
\usepackage{xcolor}
\usepackage[letterpaper, margin = 2.5cm]{geometry}
\usepackage[T2A]{fontenc}
\usepackage[english, russian]{babel}
\usepackage{amsmath, amsfonts, amssymb}
\usepackage{hyperref, url}
\hypersetup{
    colorlinks=true,
    urlcolor=blue
}

\usepackage{minted}
\usepackage{parskip}

\begin{document}

\title{ML3 HW2 Report}
\date{}
\author{Nikita Zagainov \href{https://t.me/V1adych}{@V1adych}}
\maketitle

\href{https://api.wandb.ai/links/v1adych-team/ogmbjgg0}{WandB Report}

\section*{Технические детали}
\begin{itemize}
    \item Модель: ResNet18, архитектура взята из torchvision
    \item Оптимизатор: Muon из \href{https://kellerjordan.github.io/posts/muon/}{этого поста} - реализовал руками, вдохновлялся \href{https://github.com/nil0x9/flash-muon/blob/main/flash_muon/muon.py}{здесь}
    \item Скедулер: OneCycleLR. \href{https://arxiv.org/pdf/1708.07120}{Мотивация здесь}, имплементацию брал из torch: \href{https://docs.pytorch.org/docs/stable/generated/torch.optim.lr_scheduler.OneCycleLR.html}{OneCycleLR}
    \item Логирование: логировал метрики локально в TensorBoard. Для репорта загрузил логи в \href{https://api.wandb.ai/links/v1adych-team/ogmbjgg0}{WandB}
    \item Регуляризации: применял аугментации к тренировочным картинкам, поставил \texttt{weight\_decay=0.01} в оптимизатор 
\end{itemize}

\section*{Лог экспериментов}
Единственный эксперимент - обучение модели с вышеприведенными технологиями на 100 эпохах. Не пробовал перебирать гиперпараметры или архитектурные составляющие

\section*{Результат}
Public Score в соревновании на Kaggle: 0.5255

Полученный опыт: Раньше в небольших экспериментах не использовал никаких инструментов логирования, кроме как \texttt{tqdm} или \texttt{matplotlib}, благодаря этой домашке осознал, насколько просто использовать \texttt{tensorboard} и отныне буду чаще его внедрять в свои эксперименты


\end{document}
